\chapter{The Symbolic Frontend: What You Write}\label{r01:ch:symbolic-frontend}

A \ps{} program describes a kernel with ordinary SymPy objects plus a few extensions
that \ps{} adds: \term{fields}, \term{field accesses}, \term{typed symbols} and a
handful of special functions. This chapter covers everything that can appear on the
input side of \code{ps.create_kernel}.

\section{Fields}\label{r01:sec:fields}

A \cls{Field} (\file{src/pystencils/field.py}) is the symbolic description of an array.
It has
\begin{itemize}
  \item a \textbf{name}, which becomes the keyword argument at call time;
  \item a number of \textbf{spatial dimensions} (1--3) that are iterated over;
  \item optional \textbf{index dimensions} with a fixed \textbf{index shape} (a vector or
        tensor stored in every cell, e.g.\ 19 populations of a D3Q19 lattice Boltzmann
        model);
  \item an \textbf{element type} (\code{float64}, \code{float32}, \code{int32}, a
        struct, \dots), or a \emph{dynamic} type that is resolved at code generation
        time;
  \item a \textbf{shape} and \textbf{strides}, which are either symbolic (sizes become
        kernel parameters) or fixed integers;
  \item a \textbf{memory layout} that defines which coordinate is contiguous in memory
        (\cref{r01:sec:layouts});
  \item a \textbf{field type}: \code{GENERIC} (the default), \code{INDEXED} (sparse
        index lists, \cref{r01:sec:sparse}), \code{BUFFER} (packed communication buffers),
        \code{CUSTOM} (allows absolute accesses) or \code{STAGGERED} /
        \code{STAGGERED_FLUX} (values on cell faces).
\end{itemize}

\subsection{Declaring fields}

The most convenient constructor is \code{ps.fields}, which parses a description string
\begin{center}
  \texttt{"name(index-shape), name2, ...: <data-type>[<dimensions-or-shape>]"}
\end{center}
The data type is anything \code{ps.create_type} accepts (\code{double},
\code{float32}, \code{int64}, \dots), or \code{dyn} / \code{dynidx} for the dynamic
numeric and index type. If the type is omitted (\code{"f: [2D]"}), the field gets the
dynamic numeric type. The bracket contains either a dimensionality (\code{2D}) or a
fixed shape (\code{10,20}):
\pysnippet{ex02_fields.declare}
\outsnippet{ex02_fields.declare-out}
The symbolic sizes and strides (\code{_size_pdfs_0}, \code{_stride_pdfs_0}, \dots) are
exactly the names that later appear as kernel parameters. The fixed index shape
(\code{19}) is a constant.

Fields can also be derived from existing NumPy arrays, which fixes shape, strides and
type:
\pysnippet{ex02_fields.numpy}
\outsnippet{ex02_fields.numpy-out}
The full constructors are \code{Field.create_generic(name, spatial_dimensions, dtype,
index_dimensions, layout, index_shape, field_type)},
\code{Field.create_from_numpy_array(name, array, index_dimensions, field_type)} and
\code{Field.create_fixed_size(name, shape, index_dimensions, dtype, layout, ...)}.

\subsection{Field accesses}\label{r01:sec:field-access}


\begin{figure}[htbp]
  \centering
  % Relative field accesses: offsets, compass names, and index dimensions.
\begin{tikzpicture}[font=\footnotesize]
  % ---- (a) scalar field u: the 3x3 neighbourhood -------------------------------
  \begin{scope}[x=1.55cm,y=1.15cm]
    \foreach \dx/\dy/\nm in {-1/1/NW,0/1/N,1/1/NE,-1/0/W,1/0/E,-1/-1/SW,0/-1/S,1/-1/SE} {
      \fill[psorange!25] (\dx+1,\dy+1) rectangle ++(1,1);
      \node[align=center,font=\scriptsize] at (\dx+1.5,\dy+1.5)
        {\texttt{u[\dx,\dy]}\\[-1pt]\textcolor{psgray}{\nm}};
    }
    \fill[psblue!55] (1,1) rectangle ++(1,1);
    \node[align=center,font=\scriptsize,text=white] at (1.5,1.5)
      {\texttt{u[0,0]}\\[-1pt]\texttt{u.center}};
    \draw[gridline] (0,0) grid (3,3);
    \draw[axisarrow] (0,-0.25) -- (1,-0.25) node[right,celltext] {$x$ / first offset};
    \draw[axisarrow] (-0.12,0) -- (-0.12,1) node[above,celltext] {$y$};
    \node[font=\small] at (1.5,3.35) {(a) \texttt{u = ps.fields("u: double[2D]")}};
  \end{scope}

  % ---- (b) vector field v(2): one value per index per cell -----------------------
  \begin{scope}[shift={(6.6cm,0.35cm)},x=1.55cm,y=0.75cm]
    \foreach \i in {0,1,2} {
      \draw[gridline] (\i,0) rectangle ++(1,2);
      \draw[gridline] (\i,1) -- ++(1,0);
    }
    \fill[psblue!25] (1,1) rectangle ++(1,1);
    \fill[nbcell] (2,1) rectangle ++(1,1);
    \foreach \i in {0,1,2} {
      \node[celltext] at (\i+0.5,0.5) {\texttt{(0)}};
      \node[celltext] at (\i+0.5,1.5) {\texttt{(1)}};
      \draw[domainline] (\i,0) rectangle ++(1,2);
    }
    \node[celltext] at (0.5,-0.3) {W};
    \node[celltext] at (1.5,-0.3) {C};
    \node[celltext] at (2.5,-0.3) {E};
    \draw[psgray] (1.5,2) -- (1.5,2.6) node[above,celltext] {\texttt{v(1)} = \texttt{v.center(1)}};
    \draw[psgray] (2.5,2) -- (2.9,3.3) node[above,celltext] {\texttt{v[1,0](1)}};
    \draw[decorate,decoration={brace,amplitude=3pt}] (-0.08,0) -- (-0.08,2)
      node[midway,left=4pt,align=right,celltext] {index\\dimension\\(size 2)};
    \node[font=\small] at (1.5,4.67) {(b) \texttt{v = ps.fields("v(2): double[2D]")}};
  \end{scope}
\end{tikzpicture}

  \caption[Relative field accesses]{Relative field accesses. (a) The offsets in square brackets are relative to the cell being updated; the first offset is $x$, the second $y$. (b) Round brackets select an entry of the index dimensions, here one of the two components stored in every cell.}
  \label{r01:fig:field-access}
\end{figure}

Indexing a field yields a \cls{Field.Access}, which is a subclass of
\cls{TypedSymbol} and therefore a SymPy symbol. Square brackets take \emph{relative
spatial offsets} from the current cell. Round brackets select an entry of the index
dimensions (\cref{r01:fig:field-access}):
\pysnippet{ex02_fields.access}
\outsnippet{ex02_fields.access-out}
For the centre of a scalar field, \code{u.center}, \code{u.center()}, \code{u()} and
\code{u[0,0,0]} are all the same access. The number of offsets must match the number of
spatial dimensions. Compass names
(\code{N}, \code{S}, \code{E}, \code{W}, \code{T}, \code{B} and combinations such as
\code{NE}) are short for unit offsets. Every access knows how many ghost layers it needs
(\code{required_ghost_layers}). This is how \ps{} infers the iteration space
automatically (\cref{r01:sec:iteration-space}). Offsets may also be symbolic expressions,
which gives indirect addressing such as \code{f[idx[0,0], 0]}.

\begin{note}[Absolute and staggered accesses]
  \code{f.absolute_access(offset, index)} addresses an absolute position and is only
  allowed on \code{CUSTOM} fields. Staggered fields store values on cell faces and are
  accessed with \code{f.staggered_access("E")}. The dedicated
  \code{ps.create_staggered_kernel} of \ps{}~1.x is \emph{not yet available} in 2.0
  (\cref{r01:app:pitfalls}).
\end{note}

\subsection{Memory layouts}\label{r01:sec:layouts}

The \term{layout} of a field fixes the order of its coordinates in memory, from the
slowest to the fastest varying. It determines both the strides of fixed-size fields and
the \emph{loop order} of the generated kernel: the innermost loop always runs over the
fastest coordinate. Layouts are given as strings or as tuples of dimension indices:
\begin{center}\small
\begin{tabular}{@{}lll@{}}
  \toprule
  \textbf{Name} & \textbf{Meaning} & \textbf{Fastest coordinate}\\
  \midrule
  \code{"numpy"}, \code{"c"} & C order: last index fastest & last spatial index (e.g.\ $y$ in 2D)\\
  \code{"fzyx"}, \code{"soa"}, \code{"f"} & structure of arrays & $x$; index dimension slowest\\
  \code{"zyxf"}, \code{"aos"} & array of structures & index dimension; then $x$\\
  \bottomrule
\end{tabular}
\end{center}
\Cref{r01:fig:layouts} shows the two most common layouts for a small field with two
values per cell.

\begin{figure}[htbp]
  \centering
  % Memory layouts of a 3 x 2 field with 2 values per cell: fzyx (SoA) vs. zyxf (AoS).
\begin{tikzpicture}[x=0.78cm,y=0.78cm,font=\footnotesize]
  \colorlet{fzero}{psblue!30}
  \colorlet{fone}{psorange!45}
  % ---- the field ---------------------------------------------------------------
  \begin{scope}[shift={(2.5,0.9)}]
    \foreach \x in {0,1,2} \foreach \y in {0,1} {
      \fill[fzero] (2*\x,1.2*\y) rectangle ++(1,1);
      \fill[fone] (2*\x+1,1.2*\y) rectangle ++(1,1);
      \draw[gridline] (2*\x,1.2*\y) rectangle ++(2,1);
      \draw[gridline] (2*\x+1,1.2*\y) -- ++(0,1);
      \node[celltext] at (2*\x+1,1.2*\y+1.2) {};
      \node[font=\tiny] at (2*\x+0.5,1.2*\y+0.5) {$f_0$};
      \node[font=\tiny] at (2*\x+1.5,1.2*\y+0.5) {$f_1$};
      \node[font=\tiny,text=psgray] at (2*\x+1,1.2*\y-0.12) {(\x,\y)};
    }
    \node[anchor=east,align=right,celltext] at (-0.3,1.1)
      {field \texttt{f(2)}:\\3 $\times$ 2 cells,\\2 values each};
    \draw[axisarrow] (0,-0.55) -- (2,-0.55) node[right,celltext] {$x$};
    \draw[axisarrow] (-0.15,-0.2) -- (-0.15,1.0) node[above,celltext] {$y$};
  \end{scope}

  % ---- memory strips -----------------------------------------------------------
  \newcommand{\memslot}[4]{% index, colour, label
    \fill[#2] (#1*1.0+#4,0) rectangle ++(1.0,0.8);
    \draw[gridline] (#1*1.0+#4,0) rectangle ++(1.0,0.8);
    \node[font=\tiny] at (#1*1.0+#4+0.5,0.4) {#3};}
  \begin{scope}[shift={(-1.5,-1.2)}]
    \node[anchor=east,celltext,align=right] at (-0.2,0.4) {\texttt{"fzyx"} (SoA)};
    \foreach \x/\y/\i in {0/0/0,1/0/1,2/0/2,0/1/3,1/1/4,2/1/5} {
      \memslot{\i}{fzero}{$f_0$\,\x,\y}{0}
      \memslot{\i}{fone}{$f_1$\,\x,\y}{6}
    }
    \node[anchor=west,celltext] at (12.2,0.4) {strides $x{=}1,\ y{=}3,\ f{=}6$};
  \end{scope}
  \begin{scope}[shift={(-1.5,-2.4)}]
    \node[anchor=east,celltext,align=right] at (-0.2,0.4) {\texttt{"zyxf"} (AoS)};
    \foreach \x/\y/\i in {0/0/0,1/0/1,2/0/2,0/1/3,1/1/4,2/1/5} {
      \memslot{\i}{fzero}{$f_0$\,\x,\y}{\i}
      \memslot{\i}{fone}{$f_1$\,\x,\y}{\i+1}
    }
    \node[anchor=west,celltext] at (12.2,0.4) {strides $f{=}1,\ x{=}2,\ y{=}6$};
    \draw[axisarrow] (0,-0.35) -- (12,-0.35) node[midway,below,celltext] {memory address};
  \end{scope}
\end{tikzpicture}

  \caption[Memory layouts \texttt{fzyx} and \texttt{zyxf}]{The same field \texttt{f(2)} of $3\times2$ cells in two memory layouts. With \texttt{fzyx} all values of one component are contiguous, and $x$ is the fastest coordinate: neighbouring cells are neighbours in memory, which SIMD vectorization needs. With \texttt{zyxf} the components of one cell are stored together.}
  \label{r01:fig:layouts}
\end{figure}

For a fixed-size field with shape $(4,5,6)$ and 19 values per cell, the strides (in
elements) are:
\pysnippet{ex02_fields.layouts}
\outsnippet{ex02_fields.layouts-out}
For \emph{generic} fields, whose strides are symbolic, the layout only affects the loop
order. This is why the Jacobi kernel of \cref{r01:lst:first-kernel} (default C layout)
iterates \code{ctr_1} in the inner loop, while the lattice Boltzmann kernel of
\cref{r01:ch:worked-example} (\code{fzyx}) iterates \code{ctr_0} innermost. The helper
\code{pystencils.field.create_numpy_array_with_layout(shape, layout)} allocates NumPy
arrays with a matching memory order.

\section{Assignments and assignment collections}\label{r01:sec:assignments}

\code{ps.Assignment(lhs, rhs)} is SymPy's assignment node. \ps{} extends it so that it
prints as \texttt{lhs}~$\leftarrow$~\texttt{rhs}. The left-hand side is a field access or a plain symbol. A
kernel is a list of assignments or an \cls{AssignmentCollection}, which separates
\begin{itemize}
  \item \textbf{subexpressions}: assignments to temporary symbols, emitted as local
        \code{const} variables inside the loop body, and
  \item \textbf{main assignments}: the results, usually stores to fields.
\end{itemize}
The following example writes an explicit heat-equation step with a Laplacian
subexpression \code{lap}. The diffusion coefficient \code{alpha} is a free symbol:
\pysnippet{ex03_symbolic.collection}
\outsnippet{ex03_symbolic.collection-ac}
\csnippet[nolol=false,caption={An assignment collection with a subexpression and a free scalar symbol.},label=r01:lst:heat-code]{ex03_symbolic.collection-code}
Each \emph{free} symbol, i.e.\ a symbol that is read but never assigned, becomes a
scalar kernel parameter. Its value is passed by keyword at call time:
\pysnippet{ex03_symbolic.collection-run}

\paragraph{Rules for the assignment list.}
\ps{} checks the input before generating code (\cref{r01:sec:constraints}):
\begin{itemize}
  \item every symbol is assigned at most once (\emph{static single assignment}, SSA), and
        it is not assigned after it has been read;
  \item every field location is written at most once;
  \item a field that is written at offset $o$ may be read afterwards only at the same
        offset $o$. Reading it at other offsets would make the result depend on the
        loop order.
\end{itemize}
The collection is topologically sorted, so subexpressions may be listed in any order.
Useful methods include \code{all_assignments}, \code{free_symbols},
\code{bound_symbols}, \code{new_with_substitutions(...)},
\code{new_without_subexpressions()} and \code{operation_count}.

\section{Symbols, types and casts}\label{r01:sec:types-frontend}

\ps{} is \emph{statically and strictly typed}:
\begin{itemize}
  \item A plain SymPy symbol has no type. If it is free, it gets the kernel's
        \code{default_dtype} (\code{float64} unless configured otherwise). If it is
        assigned, it gets the type of its right-hand side.
  \item \code{ps.TypedSymbol(name, dtype)} fixes the type explicitly. The type may also
        be a \code{ps.DynamicType}: \code{NUMERIC_TYPE} resolves to the default type,
        \code{INDEX_TYPE} to the index type (\code{int64} by default).
  \item Field accesses have the element type of their field.
  \item There are \textbf{no implicit conversions}. An expression mixing
        \code{float32} and \code{float64}, or an integer counter and a float, is
        rejected with a \code{TypificationError}. Use \code{ps.tcast(expr, dtype)} for
        an explicit cast.
\end{itemize}
\pysnippet{ex03_symbolic.types}
\csnippet{ex03_symbolic.types-code}
Untyped number literals take the type their context requires (\code{1.0f} in a
\code{float} context, \code{1LL} in an \code{int64} context). Types are written as strings
(\code{"float32"}, \code{"double"}, \code{"int64_t"}, \code{"uint8"}, \code{"half"}),
NumPy dtypes, Python types (\code{int} is \code{int64}, \code{float} is
\code{float64}), or \code{PsType} objects (\cref{r01:sec:types}).

\section{Supported SymPy, conditionals and special functions}\label{r01:sec:sympy-subset}

The code generator understands a well-defined subset of SymPy:
\begin{center}\small
\begin{tabularx}{\linewidth}{@{}l>{\raggedright\arraybackslash}X@{}}
  \toprule
  \textbf{Category} & \textbf{Supported}\\
  \midrule
  Numbers & \code{Integer}, \code{Rational} (becomes a division), \code{Float} (rounded to a Python \code{float}, then typed by its context), \code{pi}, \code{E}, $\pm$\code{oo}\\
  Arithmetic & \code{Add}, \code{Mul}, \code{Pow}: small integer powers ($|n|\le 8$) are expanded into multiplications, negative powers into divisions, half-integer powers use \code{sqrt}\\
  Functions & \code{exp}, \code{log}, trigonometric, hyperbolic and inverse functions, \code{atan2}, \code{Abs}, \code{floor}, \code{ceiling}, \code{Min}, \code{Max}\\
  Logic & relations (\code{<}, \code{==}, \dots) in boolean contexts, \code{And}, \code{Or}, \code{Not}, \code{true}, \code{false}\\
  Conditionals & \code{Piecewise} with a final default case \code{(expr, True)}\\
  Integers & \code{int_div}, \code{int_rem} (C semantics), bit shifts and bitwise operations from \code{pystencils.sympyextensions.integer_functions}\\
  GPU math & \code{fast_division}, \code{fast_sqrt}, \code{fast_inv_sqrt} (CUDA intrinsics)\\
  \bottomrule
\end{tabularx}
\end{center}
A \code{Piecewise} expression becomes a chain of C ternary operators. The spatial
counters \code{ps.x_}, \code{ps.y_}, \code{ps.z_} (also available as
\code{ps.DEFAULTS.spatial_counters}) give access to the cell index. Since they are
integers, they must be cast before they are combined with floats:
\pysnippet{ex03_symbolic.piecewise}
\csnippet{ex03_symbolic.piecewise-code}

\begin{pitfall}[Integer division]
  Python's \texttt{//} and \texttt{\%} round towards $-\infty$, C's \texttt{/} and \texttt{\%}
  towards zero. On integer expressions, use \code{int_div} and \code{int_rem} when you
  want the C semantics.
\end{pitfall}

\section{Reductions}\label{r01:sec:reductions}

A \term{reduction} combines a value from every cell into a single scalar. The left-hand
side must be a \code{TypedSymbol}. The operation is chosen by the assignment class
\code{ps.<Op>ReductionAssignment}, where \code{<Op>} is one of \code{Add}, \code{Sub},
\code{Mul}, \code{Min} or \code{Max}. Alternatively, call
\code{reduction_assignment(lhs, ReductionOp.<Op>, rhs)} from
\code{pystencils.sympyextensions}.
\pysnippet{ex06_reduction.reduction}
The kernel accumulates into a thread-local variable and finally combines it with the
value already stored in the result. At call time the result is therefore passed as a
\emph{one-element array}, and its initial value matters:
\pysnippet{ex06_reduction.reduction-run}
\csnippet[nolol=false,caption={Sum and maximum reductions with OpenMP.},label=r01:lst:reduction]{ex06_reduction.code}

\begin{figure}[htbp]
  \centering
  % How the OpenMP sum reduction of Listing 3.x combines all cells into total[0].
\begin{tikzpicture}[font=\footnotesize,
    acc/.style={draw=psgray!80,rounded corners=2pt,fill=white,inner sep=3pt,font=\scriptsize\ttfamily},
    res/.style={draw=psteal!80!black,fill=psteal!10,rounded corners=2pt,inner sep=4pt,font=\scriptsize\ttfamily,align=center}]
  \colorlet{th0}{psteal!45}\colorlet{th1}{psblue!35}\colorlet{th2}{psorange!45}\colorlet{th3}{psred!25}
  \begin{scope}[x=0.55cm,y=0.55cm]
    \foreach \t/\c in {0/th0,1/th1,2/th2,3/th3} {
      \fill[\c] (5*\t,0) rectangle ++(5,1);
      \node[acc,align=center] (a\t) at (5*\t+2.5,-1.7) {{\rmfamily\textcolor{psgray}{thread \t}}\\total\_local};
      \draw[flowarrow,shorten >=1pt] (5*\t+2.5,0) -- (a\t);
    }
    \draw[gridline] (0,0) grid (20,1);
    \node[anchor=east,celltext] at (-0.2,0.5) {\texttt{x}};
    \node[anchor=south west,font=\scriptsize] at (0,1.1)
      {every thread adds its share of the cells: \texttt{total\_local = total\_local + x[ctr\_0]}};
    \node[res] (omp) at (10,-4.3) {\texttt{\#pragma omp parallel reduction(+: total\_local)}\\
      combines the four partial sums};
    \foreach \t in {0,...,3} \draw[flowarrow,shorten >=1pt] (a\t.south) -- (omp);
    \node[res] (wb) at (10,-6.4) {total[0] = total[0] + total\_local};
    \draw[flowarrow] (omp) -- (wb);
    \node[anchor=west,align=left,celltext] at (16.5,-6.4)
      {\texttt{total} is a 1-element array;\\its initial value is kept};
  \end{scope}
\end{tikzpicture}

  \caption[An OpenMP sum reduction]{How the sum reduction of \cref{r01:lst:reduction} is computed: every OpenMP thread accumulates its cells into a private \texttt{total\_local}; the \texttt{reduction} clause combines the partial sums; finally the result is added to \texttt{total[0]}.}
  \label{r01:fig:reduction}
\end{figure}
\Cref{r01:fig:reduction} shows how the partial results are combined. Reductions are
supported on all CPU targets, including vectorized ones (the vector
partial results are combined horizontally, see \cref{r01:lst:rvv-reduction}), and on CUDA
through atomic operations (\cref{r01:sec:gpu-platforms}).

\section{Random numbers}

\code{ps.random.Philox(name, dtype, seed)} provides the counter-based Philox-4x32
generator~\cite{random123}. Each call to \code{get_random_vector(counter)} returns an
indexed symbol holding four \code{float32} (or two \code{float64}) values, plus the
assignment that computes them. The values depend on the counter, the cell index and the
seed. They are reproducible and need no state in memory:
\pysnippet{ex14_flow_rng.philox}
\csnippet{ex14_flow_rng.philox}
The scalar generator lives in the runtime header \file{pystencils_runtime/random.hpp}.
Vectorized variants for x86, ARM and RISC-V are in the per-ISA headers (\file{x86.hpp},
\file{neon.hpp}, \file{sve.hpp}, \file{rvv.hpp}), all built on
\file{bits/philox_rand.h}. In vectorized kernels the index type must match the width of
the random numbers, e.g.\ \code{index_dtype="int32"} for \code{float32}
(\cref{r01:sec:data-types}).

\section{From PDEs to stencils: \texttt{pystencils.fd}}\label{r01:sec:fd}

The module \code{ps.fd} discretizes partial differential equations automatically.
\code{transient}, \code{diffusion} and \code{advection} build the terms of a
conservation law. \code{Discretization2ndOrder(dx, dt)} solves for the transient term
and applies explicit Euler in time with second-order central differences in space:
\pysnippet{ex04_fd_heat.pde}
\outsnippet{ex04_fd_heat.pde}
\outsnippet{ex04_fd_heat.rhs}

\begin{figure}[htbp]
  \centering
  % The stencil produced by ps.fd for the explicit heat equation (Section 3.8).
\begin{tikzpicture}[font=\footnotesize]
  \node[align=center] (pde) at (0,0)
    {$\dfrac{\partial c}{\partial t} = D\,\nabla^2 c$\\[4pt]
     \scriptsize\texttt{transient(c) - diffusion(c, D)}};
  \node[align=center,right=0.9cm of pde] (disc)
    {\texttt{Discretization2ndOrder}\\\scriptsize explicit Euler in time,\\
     \scriptsize central differences in space};
  \draw[flowarrow] (pde) -- (disc);
  \begin{scope}[shift={($(disc.east)+(0.9,-1.275)$)},x=0.85cm,y=0.85cm]
    \foreach \dx/\dy in {1/0,-1/0,0/1,0/-1} {
      \fill[nbcell] (\dx+1,\dy+1) rectangle ++(1,1);
      \node at (\dx+1.5,\dy+1.5) {$\lambda$};
    }
    \fill[centercell] (1,1) rectangle ++(1,1);
    \node[text=white,font=\scriptsize] at (1.5,1.5) {$1{-}4\lambda$};
    \draw[gridline] (0,0) grid (3,3);
    \node[below,align=center,font=\scriptsize] at (1.5,-0.05) {$\lambda = D\,\Delta t/\Delta x^2$;\
      here $\lambda=\frac1{10}$, centre $\frac35$};
  \end{scope}
  \draw[flowarrow] (disc.east) -- ++(0.75,0);
\end{tikzpicture}

  \caption[From the heat equation to a five-point stencil]{From the heat equation to a five-point stencil: \texttt{ps.fd} discretizes the PDE, and the resulting right-hand side is a weighted sum of a cell and its four neighbours, exactly like \cref{r01:fig:stencil} with different weights.}
  \label{r01:fig:heat-stencil}
\end{figure}
With $D=1$ and $\Delta t=1/10$, the right-hand side is exactly the five-point update
$c + \tfrac{1}{10}(c_E+c_W+c_N+c_S-4c)$ (\cref{r01:fig:heat-stencil}). A time loop just swaps the arrays:
\pysnippet{ex04_fd_heat.loop}
\outsnippet{ex04_fd_heat.check}
\code{ps.fd} also provides symbolic derivatives (\code{Diff}, \code{diff}), stencil
families (\code{fd_stencils_standard}, \code{fd_stencils_isotropic}),
staggered discretizations and \code{functional_derivative} for phase-field models.

\section{Simplification and common subexpressions}\label{r01:sec:simp}

Large models, such as the lattice Boltzmann kernels in \cref{r01:ch:worked-example}, repeat
many subterms. \code{ps.simp.sympy_cse(ac)} runs SymPy's common subexpression
elimination over an assignment collection and introduces new subexpressions
\code{xi_0, xi_1, ...}:
\pysnippet{ex03_symbolic.cse}
\outsnippet{ex03_symbolic.cse-ac}
\outsnippet{ex03_symbolic.cse-counts}
The sum $\rho$ is now computed once, and the four divisions by $\rho$ became one
reciprocal \code{xi_2} and four multiplications. The backend itself does \emph{not} perform symbolic CSE. Applying it at the SymPy level
is the user's job, and it matters a lot for performance. The \code{ps.simp} module also
contains a \cls{SimplificationStrategy} class for composing rules, and helpers such as
\code{add_subexpressions_for_divisions} or \code{insert_constants}.

\section{The kernel decorator}

For short kernels, \code{@ps.kernel} offers a compact notation. The operator \code{@=}
inside the decorated function creates an assignment, and an optional argument
\code{s} creates symbols on the fly:
\pysnippet{ex03_symbolic.decorator}
\outsnippet{ex03_symbolic.decorator-asms}

\begin{pitfall}[Conditional expressions in \texttt{@ps.kernel}]
  The decorator rewrites Python's \code{a if c else b} into a \code{Piecewise}. On
  Python 3.14 this rewrite fails with \code{AttributeError: module 'ast' has no
  attribute 'NameConstant'}, because \code{ast.NameConstant} was removed from the
  standard library. Use \code{sp.Piecewise} directly instead.
\end{pitfall}

\section{Flowgraphs (experimental)}\label{r01:sec:flow}

\code{ps.flow} is a new, explicitly experimental way to write kernels as a directed
acyclic graph of \emph{equation blocks}. Inside a block, \code{_eq.let[x] = ...}
defines a private value, \code{_eq.export[y] = ...} a value visible to successor
blocks, \code{_eq.store[f(0)] = ...} a field write and \code{_eq.reduce[q, "max"] = ...}
a reduction. \code{@ps.flow.cases} creates branches. \code{ps.flow.tie} closes the graph
and checks it, and \code{create_kernel} accepts the graph directly:
\pysnippet{ex14_flow_rng.flow}
\csnippet{ex14_flow_rng.flow}
Every block keeps its own symbol table, and the generated code marks block boundaries
with comments. Flowgraphs take a separate path through the driver
(\cref{r01:sec:driver}): \code{FreezeFlowgraph} replaces \code{KernelAnalysis} and
\code{FreezeExpressions}.

\begin{keyidea}[Summary]
  \begin{itemize}[leftmargin=1.2em]
    \item A \emph{field} stands for an array; \code{f[dx, dy]} reads a neighbour relative
          to the current cell, \code{f(i)} selects a component.
    \item A kernel is a list of \emph{assignments} (or an \code{AssignmentCollection}) in
          SSA form; free symbols become scalar parameters.
    \item Types are strict: no implicit casts, use \code{TypedSymbol} and \code{tcast}.
    \item Reductions, Philox random numbers, finite differences (\code{ps.fd}) and CSE
          (\code{ps.simp}) all produce ordinary assignments.
  \end{itemize}
\end{keyidea}
